\scanpage{18}
un vectoriel \inserted{$V$} de rg $2$ sur $\mathbb{F}_2$, la donnée d'un $a\in V^*$ revient à la donnée d'une partie de $V$ qui soit une base (savoir $V^*\setminus\{a\}$); donc la donnée d'un \inserted{couple $(V,a)$ d'un} vectoriel $V$ de rg $2$ sur $\mathbb{F}_2$, et d'un élément \inserted{$a$} de $V^*$, revient à la donnée d'un ens. $C$ à deux éléments, par $V\simeq\mathbb{F}_2^C$, $a$ l'élément diagonal de $\mathbb{F}_2^C$.
Ainsi, la donnée d'un carré $Q$ \deleted{\unsure{via ses sommets}} (via l'ens des sommets) équivaut à celle d'un couple $(C,S)$, où $C$ \deleted{\unsure{est $\mathbb{F}_2$}} est un ens. à deux éléments, et $S$ est un torseur sous $\mathbb{F}_2^C$. Mais la donnée d'un tel torseur revient à celle d'une famille de $\mathbb{F}_2$-torseurs indexée par $C$, i.e. d'une application de degré $2$
\[
A\longrightarrow C,
\]
\unsure{a|à} quatre éléments, et est muni d'une involution sans pts fixes (qui donne le passage au quotient vers $C$), par la formule
\[
S\simeq\prod_{c\in C}\deleted{\mathbb{F}_2}A_c
=\text{Ens des sections de $A$ sur $C$}
\]
qui donne en même temps un plongement
\[
S\hookrightarrow\unsure{\mathcal{P}_2(A)|\mathcal{P}_2(A)_a}
\]
d'où la relation d'incidence entre $S$ et $A$\ldots
\margin{ici les deux $A_c$ ($c\in C$) \unsure{sont} \deleted{\unsure{naturellement}} \unsure{les} \uline{quotients} de $S$, par les relations d'équivalence définies resp. par \deleted{les} \inserted{les $2$ élémts} \unsure{éléments} \inserted{$c\in C$} \deleted{\unsure{deux}} de $V^*$ (considérés comme involutions ss pt fixes de $S$). Mais je préfère ici prendre la construction inverse, en \deleted{identifiant} introduisant \unsure{deux} $A_c$ et $A_{c'}$.}
% archive foliation: 17
